

\paragraph{LLM and LLM-based Autonomous Agents.} The foundation of modern LLMs is the Transformer architecture~\cite{vaswani2017attention}, which functions as the probabilistic predictor trained via the next token prediction objective. Given a sequence of tokens $x = (x_1, \dots, x_T)$, the model maximizes the joint probability:
$$
P(x) = \prod_{t=1}^{T} P(x_t \mid x_{1}, \dots, x_{t-1}; \theta).
$$
This objective enables the model to internalize world knowledge and reasoning patterns implicitly during pretraining.


While LLMs serve as the cognitive engine for reasoning and decision-making, autonomous agents extend this capability by integrating additional system components such as planning, memory, and tool-use mechanisms~\cite{wang2024survey}. These components enable agents to decompose complex tasks, retain and retrieve long-term context, and interact with external environments. In this framework, the LLM functions as the ``brain'', orchestrating actions through reasoning strategies. And agents utilize tools to perform actions beyond the model's internal knowledge.


\paragraph{Kernel Programming and Code Generation.} 
A kernel is the fundamental unit of GPU execution that concretizes high-level algorithmic semantics into hardware-level parallel operations. Since CUDA made kernels explicitly programmable, GPUs have evolved into general-purpose computing platforms, supported by highly optimized libraries such as CUTLASS~\cite{nvidia2017cutlass}. Nevertheless, writing high-performance custom kernels remains challenging, as it requires expert knowledge of hardware-specific optimization strategies. While higher-level abstractions—such as Triton and tile-based compiler frameworks—significantly improve programmability, achieving competitive performance still demands substantial domain expertise and often fails to transfer across heterogeneous accelerator platforms, underscoring the persistent gap between programmability and performance portability.

% The kernel is the fundamental unit of execution that maps high-level algorithmic semantics onto hardware. Since CUDA made kernels programmable, GPUs have evolved into general-purpose computing platforms supported by highly optimized libraries such as CUTLASS~\cite{nvidia2017cutlass}. However, writing high-performance custom kernels remains difficult, requiring expert knowledge of hardware-specific optimization. Although higher-level abstractions such as Triton and tile-based compiler frameworks improve productivity, effective performance tuning still demands substantial domain expertise and often transfers poorly across heterogeneous accelerator platforms, highlighting the persistent gap between programmability and performance portability.

In parallel, large language models have significantly advanced code generation, evolving from simple completion to managing complex software engineering workflows. However, kernel generation fundamentally differs from general-purpose code synthesis. While conventional code generation focuses on functional correctness, kernel generation must satisfy strict efficiency constraints and adapt to hardware execution characteristics. As a result, kernel generation aligns more closely with performance-oriented program synthesis and compiler optimization than with standard software development, necessitating specialized generation methods beyond generic LLM-based code generation.

% \subsection{Kernel Programming}

% As the fundamental unit of execution, the kernel maps high-level algorithmic semantics onto hardware substrates, explicitly governing workload decomposition across parallel execution units and orchestrating data movement within the memory hierarchy to saturate throughput. Following the advent of CUDA, kernels became explicitly programmable and tunable, enabling GPUs to evolve from fixed-function accelerators into general-purpose parallel computing platforms. This programmability gave rise to a rich ecosystem of highly optimized kernel libraries, including cuBLAS, cuDNN, and CUTLASS~\cite{nvidia2006geforce,lindholm2008tesla,chetlur2014cudnn,nvidia2017cutlass}.  While these libraries deliver state-of-the-art performance, writing custom kernels that approach peak efficiency remains notoriously difficult, requiring deep expertise such as tiling strategies and fine-grained synchronization.

% To mitigate this complexity, higher-level kernel abstractions have emerged. Notably, Triton~\cite{tillet2019triton} introduces a Python-based, block-level programming model that allows developers to express tile-oriented computation while delegating hardware-specific mapping and optimization to the compiler~\cite{triton2023cache}. Subsequent systems including Helion, Gluon, TileLang, and NVIDIA’s CUDA Tile/cuTile further explore compiler-mediated, tile-centric designs that trade low-level control for improved productivity and scalability across evolving hardware. However, these approaches still demand substantial domain knowledge to reason about performance. This limitation becomes more pronounced in today’s heterogeneous, multi-vendor landscape, where accelerators from AMD, Intel, and emerging vendors such as Huawei Ascend and Cambricon~\cite{hennessy2019new,reuther2020survey} exhibit divergent execution semantics and memory hierarchies. Consequently, kernels optimized for one platform often transfer poorly to others, underscoring a central limitation of existing approaches: language-level portability does not imply performance portability~\cite{pennycook2019implications}.

% \subsection{AI for Code Generation}

% Code generation applies the probabilistic modeling of LLMs to programming languages, a domain comprehensively surveyed in~\cite{jiang2024survey}. Recently, the field has evolved from simple code completion to fully integrated AI-native development environments. Tools like Cursor~\cite{cursor2025} and Trae Agent~\cite{traeresearchteam2025traeagent} have redefined the developer experience by embedding autonomous ``System 2'' reasoning directly into the IDE workflow. Furthermore, agentic command-line interfaces such as Claude code~\cite{anthropic2025claudecode} allow models to autonomously navigate file systems, execute tests, and refactor complex repositories. These advancements demonstrate a shift from generating isolated snippets to managing large-scale software engineering tasks.

% However, a critical distinction exists between general software engineering and high-performance kernel generation. General code generation prioritizes functional correctness, typically agnostic to the underlying hardware. In contrast, kernel optimization aligns closer to compiler heuristic discovery~\cite{cummins2017end} and graph-based dataflow analysis~\cite{cummins2021programl}, demanding execution efficiency and strict hardware adaptation. Consequently, generating operators requires a specialized paradigm that transcends simple correctness to optimize for latency, memory bandwidth, and micro-architectural constraints.
